% --- NUMBER SETS --------------------------------------------------------------

% inserts the symbol for the naturals
\newcommand{\N}{\mathbb{N}}

% inserts the symbol for the naturals with 0 included
\newcommand{\Nzero}{\mathbb{N}_0}

% inserts the symbol for the integers
\newcommand{\Z}{\mathbb{Z}}

% inserts the symbol for the rationals
\newcommand{\Q}{\mathbb{Q}}

% inserts the symbol for the real set
\newcommand{\R}{\mathbb{R}}

% inserts the symbol for the complex set
\newcommand{\C}{\mathbb{C}}

% inserts the symbol for the extended complex set
\newcommand{\ExtC}{\mathbb{C}_\infty}


% --- SPECIAL SYMBOLS ----------------------------------------------------------

% inserts the symbol for the identity function
\newcommand{\Id}{\mathrm{I}}


% --- DIFFERENTIATION ----------------------------------------------------------

% inserts the fraction derivative symbol, it needs two arguments:
% 1. the function to differentiante (symbol in the numerator)
% 2. the variable respect to which to differentiate (symbol in the denominator)
\newcommand{\dbyd}[2]{\frac{d#1}{d#2}}

% inserts the fraction n-th derivative symbol, it needs two arguments:
% 1. the function to differentiante (symbol in the numerator)
% 2. the variable respect to which to differentiate (symbol in the denominator)
% and can recieve an optional argument to indicate the order of the derivative
\newcommand{\dnbyd}[3][]{\frac{d^{#1}#2}{d#3^{#1}}}

% inserts the fraction derivative symbol, it needs two arguments:
% 1. the function to differentiante (symbol in the numerator)
% 2. the variable respect to which to differentiate (symbol in the denominator)
\newcommand{\pbyp}[2]{\frac{\partial#1}{\partial#2}}

% inserts the fraction n-th derivative symbol, it needs two arguments:
% 1. the function to differentiante (symbol in the numerator)
% 2. the variable respect to which to differentiate (symbol in the denominator)
% and can recieve an optional argument to indicate the order of the derivative
\newcommand{\pnbyp}[3][]{\frac{\partial^{#1}#2}{\partial#3^{#1}}}


% --- COMPLEX ANALYSIS ---------------------------------------------------------

% inserts the non italics imaginary number (roman i)
\newcommand{\I}{\mathrm{i}}

% applies the real part operator
\newcommand{\real}[1]{\mathop{Re}\left[{#1}\right]}

% applies the imaginary part operator
\newcommand{\imaginary}[1]{\mathop{Im}\left[{#1}\right]}

% --- TOPOLOGY -----------------------------------------------------------------

% inserts the diameter operator
\newcommand{\diameter}{\mathop{diam}}

% applies the interior operator
\newcommand{\interior}[1]{{#1}^{\mathrm{o}}}

% applies the closure operator
\newcommand{\closure}[1]{\overline{#1}}

% applies the boundary operator
\newcommand{\boundary}[1]{\partial {#1}}

% denotes the derived set
\newcommand{\derived}[1]{{#1}'}

% --- COMPLEX DYNAMICS ---------------------------------------------------------

% inserts the symbol for the orbit of a point
\newcommand{\orbit}{\mathcal{O}}

% applies the multiplier operator
\newcommand{\multiplier}[2]{m({#1}, {#2})}

% applies the automorphisms operator
\newcommand{\automorphism}[1]{\mathop{Aut}({#1})}

% inserts the symbol for the fixed points of a function
\newcommand{\fixedpoints}{\mathrm{Fix}}

% inserts the symbol for the periodic points of a function
\newcommand{\periodicpoints}[1]{\mathrm{Per}_{#1}}
